\documentclass[11pt]{article}
\usepackage[a4paper,margin=1in]{geometry}
\usepackage{amsmath,amssymb,amsthm,microtype}
\newtheorem{theorem}{Theorem}[section]
\newtheorem{corollary}[theorem]{Corollary}
\theoremstyle{remark}\newtheorem{remark}[theorem]{Remark}
\title{Critical Coincidence of Alias, Parity, and Gaussian Localization Scales}
\author{Bin Wang}\date{July 2026}
\begin{document}\maketitle
\begin{abstract}
When the counterloop period is tied to the intrinsic endpoint-resolution
rank, its first exact alias occurs at the same logarithmic slope as the
archived Gaussian localization ceiling.  The associated deterministic cycle
has clearance $\Theta(\sigma)$, so the old interior Gaussian replacement is
already critical.  At this order, the isolated alias bias and the scalar
square-root parity correction have the same weighted growth exponent
$0.4634069445\ldots$.  We also derive the exact first-alias matching law
necessary for a minimal-clock weighted full-trace bridge.  Separate absolute
majorants remain valid inequalities but do not decay.  No divergence of the
combined trace error is proved.
\end{abstract}

\section{Natural clock and exact alias}
Assume the typed counterloop route chooses
\[
 k_\sigma=\frac{L_\sigma}{2\log\lambda}+O(1),\qquad
 L_\sigma=\log(1/\sigma).
\]
Its exact moments are
\[
 s_{k,n}=\beta_k^n
 \bigl(2k\mathbf1_{2k\mid n}-1-(-1)^n\bigr),
 \qquad
 \beta_k\to\beta=(0.85\sqrt\lambda)^{-1}.
\]
The limiting pole moment is
$p_n=-2\mathbf1_{2\mid n}\beta^n$, and the deterministic numerator anchor
is $a_n=c_n-p_n$.

\section{Triple critical coincidence}
\begin{theorem}\label{thm:coincidence}
The first alias order satisfies
\[
 2k_\sigma=\frac{L_\sigma}{\log\lambda}+O(1).
\]
Its slope equals the archived interior-localization ceiling
\[
 a_{\rm alias}=a_{\rm loc}=\frac1{\log\lambda}
 =1.9307094191869356\ldots.
\]
The boundary cycle of physical period $2k_\sigma$ has clearance
$\Theta(\lambda^{-2k_\sigma})=\Theta(\sigma)$.  Hence its
clearance-to-noise ratio has power exponent zero, rather than tending to
infinity as required by the existing interior Gaussian proof.

At the same order, both the exact alias bias and the isolated scalar parity
correction have weighted growth exponent
\[
 \chi_1=\frac{\log(\beta R)}{\log\lambda}
 =\frac{\log(R/0.85)}{\log\lambda}-\frac12
 =0.463406944517003\ldots.
\]
\end{theorem}
\begin{proof}
The first display follows directly from the natural clock.  RH-10 and RH-17
give a boundary cycle with clearance asymptotic to a positive constant times
$\lambda^{-2k}$, proving the localization statement.

At $n=2k$,
\[
 s_{k,2k}-p_{2k}=(2k-2)\beta_k^{2k}+2\beta^{2k}.
\]
The finite-radius law
\[
 \beta_k=\beta\exp\left[-\frac{\log C_M}{2k}+o(k^{-1})\right]
\]
shows that its weighted size divided by $2k$ is
$\Theta((\beta R)^{2k})$, with exponent $\chi_1$.

For the parity eigenvalue, write
$\lambda_-(\sigma)=-(1-\delta_\sigma)$ with
$\delta_\sigma=C_*\sqrt\sigma+o(\sqrt\sigma)$.  Since
$k_\sigma\delta_\sigma\to0$,
\[
 \frac{|\lambda_-^{2k}-1|}{2k}
 \left(\frac R{0.85}\right)^{2k}
 =C_*\sqrt\sigma
 \left(\frac R{0.85}\right)^{2k}(1+o(1)).
\]
Its exponent is
$\log(R/0.85)/\log\lambda-1/2=\chi_1$.
\end{proof}

\section{Necessary first-alias matching law}
Let
\[
 e_{\sigma,k,n}=c_{\sigma,n}-s_{k,n}-a_n
 =(c_{\sigma,n}-c_n)-(s_{k,n}-p_n).
\]
The minimal bridge clock lies beyond $2k_\sigma$, so a weighted full-trace
estimate includes the first alias.

\begin{corollary}\label{cor:matching}
If the minimal-clock weighted error tends to zero, then
\[
 c_{\sigma,2k}-c_{2k}
 =(2k-2)\beta_k^{2k}+2\beta^{2k}
 +o(kR^{-2k}).
\]
The required absolute matching exponent is
\[
 \frac{\log R}{\log\lambda}
 =0.6496301165394711\ldots.
\]
\end{corollary}
\begin{proof}
The single alias term of the nonnegative weighted sum must tend to zero, so
$e_{\sigma,k,2k}=o(kR^{-2k})$.  Substitute the exact alias formula.  Since
$2k=L_\sigma/\log\lambda+O(1)$, the power of $R^{-2k}$ is the displayed
exponent.
\end{proof}

\section{Protocol and scope}
The computation independently evaluates the alias, parity, matching, and
minimal-clock clearance exponents.  The first-alias clearance exponent is
exactly zero; at the longer minimal bridge clock the clearance ratio decays
with exponent $0.4521466691\ldots$.

\begin{remark}
Triangle inequalities that majorize alias and parity pieces separately are
mathematically valid, but their bounds grow and therefore cannot close the
bridge.  A joint moving-order boundary-layer trace law could still cancel
them.  No actual full-trace divergence, spectral-submultiset identification,
or Gate A--E conclusion is proved.
\end{remark}
\end{document}
